% Created 2023-12-27 mié 22:52
% Intended LaTeX compiler: pdflatex
\documentclass[11pt]{article}
\usepackage[utf8]{inputenc}
\usepackage[T1]{fontenc}
\usepackage{graphicx}
\usepackage{longtable}
\usepackage{wrapfig}
\usepackage{rotating}
\usepackage[normalem]{ulem}
\usepackage{amsmath}
\usepackage{amssymb}
\usepackage{capt-of}
\usepackage{hyperref}
\usepackage{minted}

\usepackage{parskip}
\usepackage[utf8]{inputenc} %% For unicode chars
\usepackage{placeins}
\usepackage[margin=2.50cm]{geometry}
\usepackage[T1]{fontenc}
\usepackage{mathpazo}
\linespread{1.05}
\usepackage[scaled]{helvet}
\usepackage{courier}
\hypersetup{colorlinks=true,linkcolor=blue}
\RequirePackage{fancyvrb}
\DefineVerbatimEnvironment{verbatim}{Verbatim}{fontsize=\small,formatcom = {\color[rgb]{0.5,0,0}}}
\usepackage{mdframed}
\BeforeBeginEnvironment{minted}{\begin{mdframed}}
\AfterEndEnvironment{minted}{\end{mdframed}}
\setcounter{secnumdepth}{3}
\date{}
\title{}
\hypersetup{
 pdfauthor={},
 pdftitle={},
 pdfkeywords={},
 pdfsubject={},
 pdfcreator={Emacs 28.1 (Org mode 9.6.12)}, 
 pdflang={Spanish}}
\begin{document}

\setcounter{tocdepth}{3}
\tableofcontents

\setlength\parindent{10pt}

\section{More on Data Passing}
\label{sec:org6c62f7f}
Passing data through URLs or forms is a fundamental aspect of web
development with PHP. The two primary methods for this are GET and POST,
each with its specific use cases and characteristics.

\subsection{GET Method}
\label{sec:org3e8e0e4}
The GET method sends data through the URL. It's primarily used for
requesting data from a specific resource and is visible to the user in
the browser's address bar.

\subsubsection{Characteristics of GET:}
\label{sec:org8ee353a}
\begin{itemize}
\item Data is appended to the URL as query string parameters.
\item Suitable for non-sensitive data, as it's visible in the URL.
\item Limited amount of data can be sent (depends on the browser's URL
length limit).
\item Can be bookmarked and cached.
\item Faster compared to POST for small amounts of data.
\end{itemize}

\subsubsection{Example Usage:}
\label{sec:org3a79999}
\begin{minted}[]{html}
<!-- HTML Form using GET method -->
<form action="process.php" method="get">
    Name: <input type="text" name="name">
    <input type="submit" value="Submit">
</form>
\end{minted}

In PHP, you can access GET data using the \texttt{\$\_GET} superglobal array.

\begin{minted}[]{php}
<?php
if (isset($_GET['name'])) {
    echo "Hello, " . htmlspecialchars($_GET['name']);
}
?>
\end{minted}

\subsection{POST Method}
\label{sec:org4e18039}
The POST method sends data through the HTTP request body. It's used for
sending data to a server to create/update a resource and is not visible
in the URL.

\subsubsection{Characteristics of POST:}
\label{sec:org27c20e4}
\begin{itemize}
\item Data is not displayed in the URL.
\item Suitable for sending sensitive or large amounts of data.
\item No size limitations for the amount of data (within practical limits).
\item Cannot be bookmarked or cached.
\item Slightly slower for small amounts of data due to overhead.
\end{itemize}

\subsubsection{Example Usage:}
\label{sec:orgbb4e5ab}
\begin{minted}[]{html}
<!-- HTML Form using POST method -->
<form action="process.php" method="post">
    Name: <input type="text" name="name">
    <input type="submit" value="Submit">
</form>
\end{minted}

In PHP, POST data is accessed using the \texttt{\$\_POST} superglobal array.

\begin{minted}[]{php}
<?php
if ($_SERVER["REQUEST_METHOD"] == "POST") {
    echo "Hello, " . htmlspecialchars($_POST['name']);
}
?>
\end{minted}

\subsubsection{Conclusion}
\label{sec:org7d305a4}
\begin{itemize}
\item Use \textbf{GET} for simple, non-secure data retrieval, especially when you
want to be able to bookmark or share the URL.
\item Use \textbf{POST} for submitting forms, especially when dealing with
sensitive data or when the amount of data is considerable.
\end{itemize}

Both methods are integral to PHP web development, allowing for efficient
data transfer between client and server. Understanding when and how to
use each method is key to building robust and secure web applications.
\end{document}
